\chapter{Giới thiệu}
\label{chap:c1_introduction}

\section{Bối cảnh và lý do chọn đề tài}
\label{sec:c1_context}

Thị trường thương mại điện tử Việt Nam đạt 32 tỷ USD vào năm 2024 với tốc độ tăng trưởng 27\%~\cite{vecom2024}, trong đó mảng thiết bị công nghệ chiếm tỷ trọng lớn nhất (khoảng 35\%)~\cite{statista_vn_electronics}. Đây là mảng có đặc thù phức tạp: mỗi sản phẩm tồn tại dưới hàng chục biến thể về cấu hình, dung lượng và màu sắc. Người tiêu dùng thường có nhu cầu đa tiêu chí (ví dụ: ``laptop chạy được Premiere, pin cả ngày, dưới 20 triệu'') nhưng bộ lọc truyền thống chỉ hỗ trợ lọc theo tiêu chí rời rạc~\cite{chen2021conversational}.

Kỹ thuật Retrieval-Augmented Generation (RAG)~\cite{lewis2020rag} giải quyết khoảng trống này bằng cách kết hợp truy xuất thông tin từ cơ sở dữ liệu thực tế với khả năng sinh ngôn ngữ tự nhiên của mô hình ngôn ngữ lớn (LLM). Tuy nhiên, các nghiên cứu RAG hiện có chủ yếu tập trung vào tiếng Anh~\cite{lewis2020rag,gao2023retrieval} và chưa giải quyết đồng thời ba thách thức đặc thù của thương mại điện tử tiếng Việt: chuẩn hóa viết tắt thương hiệu (``ip'' cho iPhone, ``ss'' cho Samsung), đồng bộ vector store khi catalog thay đổi, và kết hợp Hybrid Search để tìm kiếm ngữ nghĩa và từ khóa bổ trợ lẫn nhau. Xuất phát từ khoảng trống đó, dự án này nghiên cứu và hiện thực một hệ thống thương mại điện tử hoàn chỉnh tích hợp AI chatbot theo kiến trúc RAG Hybrid cho bối cảnh tiếng Việt.

\section{Mô tả bài toán}
\label{sec:c1_problem}

Bài toán: \textit{Thiết kế và xây dựng hệ thống website thương mại điện tử chuyên bán thiết bị công nghệ, tích hợp AI Chatbot sử dụng kỹ thuật RAG để tư vấn sản phẩm theo thời gian thực, hỗ trợ thanh toán trực tuyến và cung cấp công cụ quản trị toàn diện.}

Dự án hướng đến ba câu hỏi nghiên cứu: (1) RAG Hybrid kết hợp tìm kiếm ngữ nghĩa và từ khóa có hiệu quả hơn phương pháp đơn lẻ cho tư vấn sản phẩm tiếng Việt không? (2) Bước tiền xử lý nào giảm thiểu được ảnh hưởng của viết tắt thương hiệu đến chất lượng retrieval? (3) Modular Monolith có đáp ứng được tính toàn vẹn dữ liệu và kiểm thử độc lập trong hệ thống thương mại điện tử hoàn chỉnh không?

Hệ thống phục vụ bốn nhóm tác nhân (guest, customer, staff, admin) theo nguyên tắc đặc quyền tối thiểu. Chatbot nhận câu hỏi tiếng Việt tự nhiên và trả về phản hồi kèm sản phẩm gợi ý từ catalog thực tế. Phạm vi: bán lẻ thiết bị công nghệ, giao tiếp HTTP REST, vector store đồng bộ với catalog.

\section{Mục tiêu và phạm vi nghiên cứu}
\label{sec:c1_objectives}

\subsection{Mục tiêu nghiên cứu}

Dự án đặt ra năm mục tiêu:
\begin{enumerate}
  \item Nghiên cứu và đề xuất kiến trúc RAG Hybrid cho tư vấn sản phẩm tiếng Việt, giải quyết đồng thời chuẩn hóa viết tắt, đồng bộ catalog và kết hợp tìm kiếm ngữ nghĩa với từ khóa.
  \item Thiết kế backend Modular Monolith với DI và UnitOfWork cùng frontend Feature-Based.
  \item Tích hợp cổng thanh toán MoMo và VNPay với HMAC và idempotency.
  \item Xây dựng bộ kiểm thử đa tầng ($\sim$7.303 ca kiểm thử, coverage $>$99\%).
  \item Đánh giá hiệu quả chatbot RAG qua các kịch bản thực tế đại diện.
\end{enumerate}

\subsection{Phạm vi nghiên cứu}

Về kỹ thuật, dự án tập trung vào kiến trúc Modular Monolith và Feature-Based, kỹ thuật RAG Hybrid Search, tích hợp LLM qua API, và kiểm thử đa tầng. Về chức năng, hệ thống bao gồm toàn bộ luồng mua sắm (duyệt, giỏ hàng, đặt hàng, thanh toán) và chatbot AI tư vấn sản phẩm. Hệ thống không bao gồm logistics thực tế, collaborative filtering hay quản lý kho vật lý.

\subsection{Phương pháp nghiên cứu}

Dự án áp dụng phương pháp nghiên cứu thiết kế (design science research) qua bốn giai đoạn: khảo sát lý thuyết RAG và công nghệ liên quan, thiết kế kiến trúc và pipeline RAG Hybrid, hiện thực hệ thống kèm kiểm thử đa tầng, và đánh giá hiệu quả chatbot qua kịch bản thực tế.

\section{Cấu trúc dự án}
\label{sec:c1_structure}

Dự án gồm bốn chương: \textbf{Chương 1} trình bày bối cảnh và mục tiêu; \textbf{Chương 2} tổng hợp nền tảng lý thuyết RAG, Hybrid Search và kiến trúc phần mềm; \textbf{Chương 3} phân tích yêu cầu, thiết kế kiến trúc, cơ sở dữ liệu và pipeline RAG; \textbf{Chương 4} trình bày kết quả thực nghiệm gồm cài đặt, kiểm thử và đánh giá hiệu quả chatbot. Phần \textbf{Kết luận} tổng kết kết quả và đề xuất hướng phát triển.
